% Canonical continuation, not a second explanatory pass: C01.06--C01.13.
\subsection{সফল যোগাযোগের চারটি মৌলিক গুণ}
\label{ch2:c01-06}
\chTwoPair{\textbf{সঠিক গন্তব্যে পৌঁছানো (Delivery):} ডেটা নির্ধারিত প্রাপকের কাছে পৌঁছাতে হবে। ভর্তি ফরম অন্য কলেজে পৌঁছালে আবেদনের উদ্দেশ্য পূরণ হবে না।\par
\textbf{নির্ভুলতা (Accuracy):} ডেটা বিকৃত না হয়ে প্রাপকের কাছে পৌঁছাতে হবে। সঠিক প্রাপকের কাছে গেলেও টাকার অঙ্ক বদলে গেলে লেনদেন ভুল হবে।}{\textbf{সময়মতো পৌঁছানো (Timeliness):} ডেটা প্রয়োজনীয় সময়ের মধ্যে পৌঁছাতে হবে। পরীক্ষার নির্ধারিত সময় শেষ হওয়ার পর উত্তর পৌঁছালে তা কাজে আসে না।\par
\textbf{বিলম্বের ওঠানামা (Jitter):} ধারাবাহিক প্যাকেট পৌঁছাতে যে বিলম্ব হয়, তার ওঠানামাকে জিটার বলে। জিটার বেশি হলে ভিডিও ক্লাসের শব্দ ভেঙে আসতে পারে। তাই কম জিটার কাম্য।}
সফল যোগাযোগে নিরাপত্তা ও নির্ভরযোগ্যতাও প্রয়োজন। অনুমতিহীন ব্যবহার ঠেকাতে হয় এবং প্রয়োজনের সময়ে সেবা সচল রাখতে হয়। তবে যোগাযোগের প্রচলিত চারটি মৌলিক গুণ হলো সঠিক গন্তব্যে পৌঁছানো, নির্ভুলতা, সময়মতো পৌঁছানো এবং কম জিটার। মনে রাখবে, বিলম্ব বেশি হওয়া ও বিলম্বের ওঠানামা আলাদা সমস্যা।

\subsection{সংকেত দুর্বল বা বিকৃত হওয়ার কারণ}
\label{ch2:c01-07}
\begin{tabularx}{\linewidth}{|p{.20\linewidth}|X|X|}
\hline
\textbf{সমস্যা} & \textbf{কী হয়} & \textbf{করণীয়ের উদাহরণ}\\\hline
সংকেতের ক্ষয় (Attenuation) & দূরত্ব বাড়লে সংকেতের শক্তি কমে। & উপযুক্ত মাধ্যম বা সংকেত পুনর্গঠনের ব্যবস্থা।\\\hline
সংকেতের বিকৃতি (Distortion) & সংকেতের আকৃতি বদলায়; বিভিন্ন অংশের বিলম্ব ভিন্ন হতে পারে। & উপযুক্ত মাধ্যম, সংকেতের উপযোগী বিন্যাস বা এনকোডিং (Encoding) এবং গ্রহণব্যবস্থা।\\\hline
নয়েজ (Noise) & অনাকাঙ্ক্ষিত সংকেত মূল সংকেতের সঙ্গে মেশে। & বিঘ্নের উৎস থেকে দূরে রাখা, সুরক্ষামূলক আবরণ বা শিল্ডিং (Shielding) এবং ত্রুটি শনাক্তকরণ।\\\hline
\end{tabularx}
\textbf{EMI (Electromagnetic Interference)} হলো বাইরের বিদ্যুৎচৌম্বক ক্ষেত্রের কারণে সংকেতে সৃষ্ট বিঘ্ন। মোটরের পাশে থাকা তামার কেবলে এ ধরনের সমস্যা হতে পারে। অপটিক্যাল ফাইবারের আলোকপথ EMI-তে প্রভাবিত হয় না। তবে অতিরিক্ত বাঁক, কাটা পড়া বা সংযোগস্থলের ক্ষতিতে ফাইবারের সংকেতে সমস্যা হতে পারে। কেবলের গঠন ও রিপিটারের কাজ সংশ্লিষ্ট অংশে আলোচনা করা হয়েছে।

\subsection{ব্যান্ডউইথ: ক্ষমতা ও ফ্রিকোয়েন্সির প্রস্থ}
\label{ch2:c01-08}
\begin{definitionbox}{Bandwidth}
ডিজিটাল নেটওয়ার্কের প্রচলিত আলোচনায় সংযোগের সর্বোচ্চ ডেটা বহনের হারকে ব্যান্ডউইথ বলা হয়। এই প্রসঙ্গে ব্যান্ডউইথ দিয়ে সংযোগের বহনক্ষমতা (Capacity) বোঝানো হয়। এর একক বিট প্রতি সেকেন্ড বা bits per second (bps)। ফ্রিকোয়েন্সির ক্ষেত্রে ব্যান্ডউইথ হলো সর্বোচ্চ ও সর্বনিম্ন ফ্রিকোয়েন্সির পার্থক্য; এর একক হার্টজ (Hertz, Hz)।
\end{definitionbox}
\chTwoPair{\textbf{ডেটা বহনের হার:} 10 Mbps সংযোগের নামমাত্র বহনক্ষমতা প্রতি সেকেন্ডে 10 মিলিয়ন বিট। বাস্তবে ব্যবহৃত অ্যাপ্লিকেশনে এই পুরো হার পাওয়া নাও যেতে পারে।}{\textbf{ফ্রিকোয়েন্সির প্রস্থ:} সর্বনিম্ন ফ্রিকোয়েন্সি 1000 Hz এবং সর্বোচ্চ ফ্রিকোয়েন্সি 4000 Hz হলে ব্যান্ডউইথ $4000-1000=3000$ Hz। Hz ও bps ভিন্ন পরিমাণের একক; সরাসরি সমান ধরে নেওয়া যায় না।}
বহনক্ষমতা বেশি হলে একই সময়ে বেশি ডেটা পাঠানোর সুযোগ থাকে। তবে বাস্তবে পাওয়া হার মাধ্যম ও যন্ত্রের সামর্থ্য, অন্য ব্যবহারকারীর চাপ, সার্ভারের সাড়া এবং যোগাযোগের বিলম্বের ওপরও নির্ভর করে।

\subsection{bit, Byte ও হারের একক}
\label{ch2:c01-09}
বিট (bit)-এর পূর্ণরূপ \textbf{Binary Digit}; এর মান 0 বা 1। $1\ \mathrm{Byte}=8\ \mathrm{bit}$। এককের চিহ্নে ছোট হাতের b দিয়ে বিট এবং বড় হাতের B দিয়ে বাইট বোঝানো হয়। তাই 8 Mbps ও 8 Megabyte/s (MB/s) একই হার নয়।
\begin{tabularx}{\linewidth}{|p{.18\linewidth}|X|p{.28\linewidth}|}
\hline
\textbf{চিহ্ন} & \textbf{পূর্ণরূপ} & \textbf{মান}\\\hline
bps & bits per second & বিট/সেকেন্ড\\\hline
kbps (Kbps) & kilobits per second & $10^3$ bps\\\hline
Mbps & megabits per second & $10^6$ bps\\\hline
Gbps & gigabits per second & $10^9$ bps\\\hline
Tbps / Pbps & terabits / petabits per second & $10^{12}$ / $10^{15}$ bps\\\hline
\end{tabularx}
নেটওয়ার্কে ডেটা বহনের হার প্রকাশে দশমিক একক ব্যবহৃত হয়। দশমিক হিসাবে MB হলো $10^6$ Byte; অন্যদিকে, মেবিবাইট বা Mebibyte (MiB) হলো $2^{20}$ Byte। প্রশ্নে MB-কে $1024^2$ Byte ধরা হলে সেই প্রদত্ত শর্ত অনুযায়ী হিসাব করতে হবে। দশমিক হারের সঙ্গে বাইনারি ভিত্তিতে প্রকাশিত ফাইলের আকার ব্যবহার করার আগে এককের ভিত্তি মিলিয়ে নিতে হবে।

\subsection{প্রচলিত ব্যান্ডউইথ শ্রেণি}
\label{ch2:c01-10}
নিচের ছকে পরীক্ষায় প্রচলিত ব্যান্ডউইথের শ্রেণিবিন্যাস দেখানো হয়েছে। আধুনিক প্রযুক্তিতে ব্রডব্যান্ডের সংজ্ঞা ও গতি মানদণ্ডভেদে ভিন্ন হতে পারে।
\begin{tabularx}{\linewidth}{|p{.22\linewidth}|p{.25\linewidth}|X|}
\hline
\textbf{শ্রেণি} & \textbf{প্রচলিত হার} & \textbf{ব্যবহার ও পরিচয়}\\\hline
Narrowband & 45--300 bps & টেলিগ্রাফ ও ধীর বার্তা আদান-প্রদান।\\\hline
Voiceband & 1200--9600 bps & টেলিফোন-চ্যানেলে প্রচলিত ডেটা যোগাযোগ।\\\hline
Broadband & 1 Mbps বা বেশি & বেশি ডেটা বহন: কেবল, ফাইবার ও উপযুক্ত বেতার সংযোগ।\\\hline
\end{tabularx}
bps থেকে kbps-এ রূপান্তর করতে 1000 দিয়ে ভাগ করতে হয়। যেমন, $9600$ bps $=9.6$ kbps এবং $300$ bps $=0.3$ kbps। মনে রাখবে, এই ঐতিহাসিক শ্রেণিবিন্যাসে দেওয়া 9600 bps টেলিফোনের সব প্রযুক্তির সর্বোচ্চ গতিসীমা নয়।
\textbf{ওয়াইডব্যান্ড (Wideband)} দিয়ে সাধারণত বিস্তৃত ফ্রিকোয়েন্সি-পরিসরের সংকেত বা চ্যানেল বোঝানো হয়। তাই ওয়াইডব্যান্ডকে ওপরের তিনটি শ্রেণির মতো নির্দিষ্ট ডেটা হারের চতুর্থ শ্রেণি বলা যায় না। Hz-এ ব্যান্ডউইথ নির্ণয় করতে প্রশ্নে দেওয়া সর্বোচ্চ ও সর্বনিম্ন ফ্রিকোয়েন্সির পার্থক্য নিতে হবে।

\subsection{Capacity, Throughput, Latency ও Bottleneck}
\label{ch2:c01-11}
\begin{tabularx}{\linewidth}{|p{.23\linewidth}|X|}
\hline
\textbf{ধারণা} & \textbf{অর্থ ও উদাহরণ}\\\hline
বহনক্ষমতা (Capacity) & সংযোগের সর্বোচ্চ ডেটা বহনের হার; যেমন নামমাত্র 10 Mbps।\\\hline
থ্রুপুট (Throughput) & নির্দিষ্ট সময়ে সফলভাবে পৌঁছানো ডেটার হার; বাস্তবে 6 Mbps হতে পারে।\\\hline
বিলম্ব (Latency) & প্রেরক থেকে প্রাপকের কাছে ডেটা পৌঁছাতে যে সময় লাগে; কম ডেটা হারের সমস্যা থেকে এটি আলাদা।\\\hline
বটলনেক (Bottleneck) & পথে সবচেয়ে কম বহনক্ষমতার অংশ, যা পুরো পথের ডেটা হার সীমিত করে।\\\hline
\end{tabularx}
ডেটার সঙ্গে পাঠানো নিয়ন্ত্রণতথ্য বা হেডার (Header), অন্য ব্যবহারকারীর ডেটা চলাচল, পুনঃপ্রেরণ এবং সার্ভারের সীমাবদ্ধতা অ্যাপ্লিকেশনে পাওয়া হার কমায়। বহনক্ষমতা ও বিলম্বের পার্থক্য একটি রাস্তার উদাহরণে বোঝা যায়। চওড়া রাস্তা দিয়ে বেশি গাড়ি চলতে পারে, কিন্তু যানজট হলে গন্তব্যে পৌঁছাতে দেরি হয়।
\chTwoPair{\textbf{ভাগ করা সংযোগ (Shared Connection):} একাধিক গ্রাহক সংযোগের একই অংশের বহনক্ষমতা ভাগ করে ব্যবহার করেন। ২০ Mbps-এর একটি সংযোগ চারজন সমানভাবে ব্যবহার করলে সরল হিসাবে জনপ্রতি ৫ Mbps পাওয়া যায়। বাস্তবে চাহিদা ও নিয়ন্ত্রণব্যবস্থা অনুযায়ী এই ভাগ পরিবর্তিত হতে পারে।}{\textbf{সংরক্ষিত সংযোগ (Dedicated Connection):} চুক্তিতে নির্ধারিত বহনক্ষমতা গ্রাহকের জন্য সংরক্ষিত থাকে। তবে নিজের সংযোগে এই হার সংরক্ষিত থাকলেও পুরো পথের অন্য অংশ, সার্ভার বা প্রোটোকলের সীমাবদ্ধতায় ফাইল ডাউনলোডের প্রকৃত হার কমতে পারে।}

\subsection{সেবাভেদে প্রচলিত Data Rate-এর উদাহরণ}
\label{ch2:c01-12}
নিচের ছকে বিভিন্ন সেবার প্রচলিত ডেটা হারের উদাহরণ দেওয়া হয়েছে। নির্দিষ্ট মানভিত্তিক প্রশ্নে সেবার নামের সঙ্গে হারটি মিলিয়ে দেখতে হবে। বাস্তবে প্রয়োজনীয় হার সফটওয়্যার, সেবার মান ও ব্যবহারের ধরন অনুযায়ী পরিবর্তিত হয়। তাই ছকের মানগুলো নিশ্চিত সর্বনিম্ন চাহিদা নয়।
\begin{tabularx}{\linewidth}{|p{.25\linewidth}|X|}
\hline
\textbf{উদাহরণ হার} & \textbf{পরিষেবা}\\\hline
0.5 Mbps & VoIP ফোন কল, ই-মেইল ও গান স্ট্রিমিং।\\\hline
0.5--1 Mbps & ওয়েব ব্রাউজিং।\\\hline
0.7 Mbps & ভিডিও স্ট্রিমিং।\\\hline
1 Mbps & সাধারণ ভিডিও কনফারেন্সিং ও ইন্টারনেট গেম কনসোল।\\\hline
1.5 Mbps & চলচ্চিত্র স্ট্রিমিং।\\\hline
4 Mbps & HD চলচ্চিত্র স্ট্রিমিং, HD ভিডিও কনফারেন্সিং ও অনলাইন HD মাল্টিপ্লেয়ার গেম।\\\hline
\end{tabularx}
\textbf{VoIP} = Voice over Internet Protocol; \textbf{HD} = High Definition। \textbf{0.7 Mbps ভিডিও স্ট্রিমিংয়ের উদাহরণ হার}; এই ছকে VoIP-এর জন্য অন্য মান দেওয়া আছে। বাস্তবে অডিও–ভিডিও সংকোচন ও উদ্ধারের ব্যবস্থা বা কোডেক (Codec), ছবির বিস্তারিত প্রকাশের মাত্রা বা রেজল্যুশন (Resolution), ডেটা সংকোচন (Compression) এবং একসঙ্গে ব্যবহারকারীর সংখ্যা অনুযায়ী প্রয়োজনীয় হার বদলায়। অনলাইন ক্লাসের সংযোগ বেছে নেওয়ার সময় তাই শুধু এই ছক যথেষ্ট নয়। আপলোডের হার ও যোগাযোগের বিলম্বও দেখতে হয়।

\subsection{Baud ও Bit Rate-এর পার্থক্য}
\label{ch2:c01-13}
\begin{definitionbox}{Baud (Bd)}
প্রতি সেকেন্ডে পাঠানো সংকেত-প্রতীকের হার প্রকাশের একক বড (Baud)। একটি সংকেত-প্রতীক বা সিম্বল (Symbol) এক বা একাধিক বিট প্রকাশ করতে পারে। তাই বড ও bps-এর মান সবসময় সমান হয় না।
\end{definitionbox}
প্রতিটি সিম্বলে সমান সংখ্যক বিট প্রকাশিত হলে,
\[
\text{Bit Rate}=\text{Symbol Rate}\times\text{bits per symbol}.
\]
2400 Baud-এ প্রতিটি সিম্বল দুই বিট প্রকাশ করলে ডেটার হার $2400\times2=4800$ bps। প্রতিটি সিম্বল এক বিট প্রকাশ করলে বড ও bps-এর সংখ্যাগত মান সমান হয়। মনে রাখবে, বড ও বিট হার নির্ণয়ের এই হিসাব ক্যারেক্টারের ডেটা, স্টার্ট ও স্টপ বিট গণনার বিকল্প নয়। ফ্রেমিং ও কার্যকর ডেটার অনুপাত পরের অংশে আলোচনা করা হয়েছে।
